% !TeX root = ../Book4.tex
\begin{figure}[htbp]
	\centering
	\begin{tikzpicture}[font=\small,scale=0.9]
		\fill[GeometryGrayLight] (0,0) circle (2.6);
		\fill[GeometryBlueLight] (0,0) circle (1);
		\draw[GeometryGray,thick] (0,0) circle (2.6);
		\draw[GeometryGray,dashed] (0,0) circle (2);
		\draw[GeometryBlue] (0,0) circle (1);
		\fill (0,0) circle (1.6pt);
		\node[below left] at (0,0) {\(0\)};
		\node[below] at (0,-1) {\(B_r\)};
		\node[below] at (0,-2) {\(B_{2r}\)};
		\node[below] at (0,-2.6) {\(B_R\)};
		\fill[GeometryBlue] (0.65,0.35) circle (2pt);
		\node[left] at (0.6,0.4) {\(x_0\)};
		\draw[GeometryBlue,thick,->] (0.65,0.35)
			.. controls (0.75,0.5) and (0.85,0.58) .. (1,0.55);
		\node[above right] at (1,0.55) {\(\phi_t(x_0)\)};
	\end{tikzpicture}
	\caption[Darboux 构造中的局部流范围]{在 \(B_R\) 上构造向量场，从更小的 \(B_r\) 出发。二次速度估计保证轨道在整个单位时间内留在 \(B_{2r}\) 中。图为嵌套 Euclid 球及一条轨道的二维示意，曲线形状不代表特定方程的解。}
	\label{g4:fig:0802:flow-domains}
\end{figure}
